\section{The compact harmonic limit}\label{sec:global-completion}

The local estimates identify the geometry near the
exceptional cycles. We finish by identifying the
completion of the regular harmonic metric and the
limit of the entire compact spaces. In the isolated
example a direct collapsing-neighborhood argument
suffices. For the interacting family we identify
the limiting semiample class and apply the
Ricci-flat degeneration theorems of Tosatti and
Song to the actual harmonic metrics.

\subsection{Smooth convergence on the stable open set}
Let $U=Z_{\rm reg}$, identified with the
$H$-stable bundle locus, and write $g_0$ for its
central harmonic metric. A compact subset of
$U$ is covered by finitely many local universal
families. At each central bundle the normalized
fixed-determinant HE equation has invertible
trace-free gauge linearization, since stability
removes its kernel. Compactness supplies a
uniform positive lower bound. The parameter
implicit-function theorem and elliptic
bootstrapping give smooth variation of the
normalized HE metrics. The deformation kernels
have constant dimension on this smooth locus,
so their harmonic projections vary smoothly
as well. Consequently
\begin{equation}\label{eq:regular-smooth-limit}
 g_\eta\longrightarrow g_0
 \quad\text{in every fixed }C^k\text{ norm on compact subsets of }U.
\end{equation}
These are the same holomorphic bundles in every
nearby chamber. The identification is the
restriction of the classifying contraction.

\subsection{A direct argument for the isolated plane}
In the isolated zero-level model, distance to
the cone vertex is the norm radius. The local
harmonic comparison shows that a radius-$\rho$
neighborhood has diameter $O(\rho)$ and that
adding the vertex is exactly the local metric
completion, with its quotient topology.
Outside such a neighborhood the target is
compact and smooth. Thus the completion of
$(U,g_0)$ is $Z$, with no further end.

One can also see compact convergence directly.
Use the invariant $A=u\otimes\varphi$ from
\eqref{eq:isolated-incidence}, and take the
target neighborhood where its nonzero singular
value is at most $\rho^2/2$. In the balanced
model, if $\kappa_t\le\rho^2$, its whole
preimage has radius $O(\rho)$ and diameter
$O(\rho)$, uniformly in $t$. Whole-fibre
coverage and the ambient comparison transfer
this bound to the actual moduli space.

Collapse each of these neighborhoods to a
point. A shortest quotient path needs to
visit that point at most once; lifting the
visit across a set of diameter $O(\rho)$
changes distance by at most $O(\rho)$.
For fixed $\rho$, the common compact exterior
has smoothly convergent metrics by
\eqref{eq:regular-smooth-limit}, so the
collapsed spaces converge in length distance.
The limsup of the original compact
Gromov--Hausdorff distance is therefore at
most $C\rho$. Letting $\rho\downarrow0$
proves convergence to the central completion.

\subsection{The limiting class gives the actual contraction}
In the interacting example put $c=D_1+D_2$
and $b_0=(0,H,0)\in v^\perp$. On every
chamber,
\[
 \frac{[\Omega_\eta]}{2\pi^2}
       \longrightarrow\theta_v(b_0).
\]
To identify this limiting class globally, use
the Greb--Toma determinant construction with
a sufficiently large very ample multiple $mH$.
Its surface $K$-class is
\[
 u_1=-3[\OO_{mH}]
       +\chi(E\otimes\OO_{mH})[\OO_x].
\]
Riemann--Roch gives
\[
 v(\OO_{mH})=(0,mH,-(mH)^2/2),\qquad
 v(u_1)=(0,-3mH,c\cdot mH).
\]
Their determinant convention is
$\lambda_{\cE}(u_1)=\det Rp_*(\cE\otimes p_X^*u_1)$
\cite[Definition~2.1]{grebtoma}. Since our
Mukai convention inserts the dual vector,
and $c\cdot H=0$, we obtain
\begin{equation}\label{eq:determinant-polarization}
 c_1\lambda_{\cE}(u_1)=3m\theta_v(b_0).
\end{equation}

The classifying theorem gives an ample line
on the slope target whose pullback is
$\lambda_{\cE}(u_1)^N$ for a positive integer
$N$. Restrict to the common reduced image
and pull back under its finite normalization.
The resulting line $A_Z$ on $Z$ is ample.
Using a common power for the finitely many
chambers gives
\begin{equation}\label{eq:common-ample-class}
 \theta_v(b_0)=\pi_\sigma^*\alpha_Z,
 \qquad \alpha_Z=\frac{c_1(A_Z)}{3mN}.
\end{equation}
The maps here are precisely the classifying
contractions already constructed. Thus the
limiting class is big and nef, and its
sufficiently divisible multiples define
this connected-fibre semiample morphism.
Vanishing on the named exceptional curves
alone would not have identified the ample
class on the target.

\subsection{Ricci-flatness and metric identification}
Every chamber object is locally free and
slope stable. At fixed determinant its
trace-free HE connection is an irreducible
projective ASD connection on the Ricci-flat
K3 surface. Itoh's construction
\cite[Theorem~2.1]{itoh} makes its smooth
moduli metric hyperk\"ahler. The real
Dolbeault metric differs by the fixed factor
specified in Section~\ref{sec:harmonic-models};
it is therefore hyperk\"ahler as well,
smooth on the entire compact $M_\sigma$,
and complete. By uniqueness in a K\"ahler
class,
$g_\eta/(2\pi^2)$ is the Ricci-flat metric
in class $\theta_v(b_\eta)$ on that fixed
complex manifold. No auxiliary metric on
a deformation-equivalent space is used.

The symplectic form on $U$ pulls back and
extends as the symplectic form on
$M_\sigma$. Its top power trivializes the
canonical sheaf on $U$. Normality and the
codimension-at-least-two singular complement
give the trivial reflexive canonical sheaf
on $Z$. The resolution has zero discrepancies;
further resolutions have nonnegative
discrepancies, since $M_\sigma$ is smooth.
Hence $Z$ is a projective Calabi--Yau
variety with canonical singularities and
a crepant resolution, as required by the
metric completion results.

Tosatti's real-class degeneration theorem
\cite[Theorem~1.1 and the subsequent sequence
formulation]{tosatti} applies to
$\theta_v(b_\eta)\to\pi_\sigma^*\alpha_Z$.
Its smooth limit away from the exceptional
locus is the regular Ricci-flat metric
of $\alpha_Z$ on $Z$. Equation
\eqref{eq:regular-smooth-limit} identifies
that metric with $g_0/(2\pi^2)$.
This equality of metric tensors is a separate
step from the algebraic/Uhlenbeck topological
identification.

\begin{proposition}\label{prop:compact-limit}
The completion of $(U,g_0)$ is a compact
length space $(Z,d_0)$ whose topology is
the projective topology. For every generic
sequence $\eta\to0$, allowing changes of
chamber and arbitrary ratios of its positive
scales,
\[
 (M_{\sigma(\eta)},g_\eta)
          \longrightarrow(Z,d_0)
 \quad\text{in compact Gromov--Hausdorff distance}.
\]
\end{proposition}

\begin{proof}
Song's completion theorem
\cite[Theorem~1.1]{song} applies to
$(Z,\alpha_Z)$. Its regular Ricci-flat
metric has compact completion homeomorphic
to $Z$, extending the identity on $U$.
The metric identification above and a fixed
homothety give the first assertion.
On a fixed chamber, \cite[Corollary~1.1]{song}
gives compact convergence of the Ricci-flat
metrics to that completion. Its limiting
class is rational, big and nef, with
semiample model exactly $Z$ by
\eqref{eq:common-ample-class}.

The conclusion applies to real interior
classes as well. If rational approximants
are used in this argument, approximate the
$n$th interior metric by a rational interior
class with Gromov--Hausdorff discrepancy
less than $1/n$. Smooth dependence at each
interior class permits this choice, and the
approximating classes still tend to the
same boundary class. No continuity estimate
uniform up to the boundary is assumed.

There are only six chambers. Every generic
parameter sequence has subsequences in a
single chamber. The determinant calculation
and the regular harmonic identification
give the same $(Z,d_0)$ for all such
subsequences. The subsequence criterion
proves the assertion for the original
sequence.
\end{proof}

\begin{proof}[Proof of Theorem~\ref{thm:main}]
Section~\ref{sec:compact-interaction} constructs
the compact chamber spaces, excludes all
ideal charge and identifies the full
singular locus of their common normal
target. Section~\ref{sec:chambers} gives
all six deepest fibres and the reduced
specialization. The primitive cycle
calculations and Proposition
\ref{prop:calculated-attachments} prove
the integral assertions. Equations
\eqref{eq:exact-cycle-volumes},
\eqref{eq:relative-plane-diameter},
\eqref{eq:uniform-plane-bounds} and
\eqref{eq:separated-diameter}, together
with Theorem~\ref{thm:successive-limits},
give the exact periods, ambient collapse
and complete rescaled limits for those
same cycles. Proposition~\ref{prop:compact-limit}
proves the final compact convergence.
\end{proof}

The integral attachment uses a marked
complement, and the metric completion uses
the global harmonic class on the actual
compact family. The exceptional cycles
thus carry explicit global information,
while the geometry retained away from
them remains part of the reconstruction
data. The argument concerns these
contractions without lost charge. Extending
it to singular graded sheaves and bubbling
would require compatible local metrics and
integral attachment data for those additional
boundary strata.
